\section{Conclusion}
\label{sec:conclusion}

We have characterised the solution space of minimum-edge-cut redistricting
for US congressional districts and drawn three conclusions relevant to the
Districting Integrity Act and to the broader question of neutral algorithmic
redistricting.

\paragraph{The seed is irrelevant for certified states.}
For every focal state in our analysis, the minimum edge-cut plan converges:
running minima improve infrequently and eventually plateau. For certified
states (stable tail $\geq 300$ seeds), the partisan outcome does not change
regardless of which specific seed achieves the minimum. The seed is a
computational shortcut to reproduce the converged plan, not a hidden partisan
choice. The correct statutory protocol is not ``run $N$ seeds'' but ``run
seeds until convergence is certified'': our focal-state sweeps reveal that
interim results at 200 seeds are overturned in three states (Georgia at seed
489, Tennessee at seed 609, Wisconsin at seed 318), underscoring that fixed
seed budgets cannot substitute for convergence testing.

\paragraph{Geometric sorting drives proportionality, not algorithmic bias.}
The minimum-edge-cut criterion produces near-proportional outcomes in
geographically mixed states (Georgia $-0.1$\,pp, Michigan $+2.4$\,pp,
Pennsylvania $-3.5$\,pp) and disproportional outcomes in sorted states
(Wisconsin $-12.8$\,pp, North Carolina $-13.6$\,pp, Tennessee $-27.1$\,pp).
This pattern is not algorithmic bias: any criterion that rewards geographic
compactness will produce the same structure, because compact districts respect
geographic sorting. The proportionality gap measures the geographic
distribution of voters, not the algorithm's fairness. Achieving proportional
representation in sorted states requires importing partisan data, which the
Districting Integrity Act expressly prohibits for federal congressional
districts (\S~104(e)). We recommend reporting the proportionality gap as a
mandatory transparency metric alongside the reproducibility artifacts, not as
a standard the algorithm is required to meet.

\paragraph{Edge-cut is the correct primary criterion.}
Among the geometric criteria studied --- minimum edge-cut, maximum
geometric-mean Polsby-Popper, and proportional bisection --- edge-cut is
uniquely defensible. It has no interpretive degrees of freedom: the inputs
(TIGER boundary lengths) are federal public records with no alternative
measures, and the objective (minimise total boundary length cut) has a single
unambiguous definition. Compactness metrics introduce choices (which area
measure? which shape metric?) that could be argued to have been selected
post-hoc. Proportional bisection requires partisan data as input, prohibited
by \S~104(e). We propose the \emph{minimum-edge-cut protocol} as the
standard: run seeds until convergence is certified, publish the minimum-EC
plan with its convergence trajectory.

\paragraph{Limitations and future work.}
The Cheeger/Fiedler lower bound is mathematically valid but empirically
non-certifying for US census-tract graphs: all states have near-zero
$\lambda_2$ (7--46\,m) due to geographic pinch-points, giving bounds 50--200×
below the achieved minimum. Tighter bounds from spectral graph theory (SDP
relaxations, higher Cheeger constants) or from the graph's isoperimetric
profile would enable formal mathematical certification complementing our
empirical convergence certificates. The \textsc{CompactBisect} algorithm
achieves the same converged partisan outcome as MEC for all tested states, via
multi-level geometric-mean PP optimisation, and merits further study as a
complementary criterion that reaches convergence faster. Finally, the 14\%
failure rate for Texas and 64\% for California under the flat 0.5\% tolerance
schedule motivates implementation of the statute's tiered tolerance
(\S~104(b)(4)) for large-state deployments.

\vspace{1em}
\noindent\textbf{Reproducibility.} All code, adjacency files, seed-sweep
outputs, and convergence data are available at the reference implementation
repository. The convergence trajectories for all focal states (seed $\times$
edge-cut tables) are published as reproducibility artifacts alongside this
paper.
